% ============================================================
% PASTE THIS INSIDE \subsection{Index Variables}
% Replace the single placeholder line:
%   "Describe the index role, the evidence supporting..."
% ============================================================

\label{sec:index}
For index and key variables, we probe Qwen2.5-1.5B and DeepSeek-Coder-1.3B on
the same XLCoST corpus using the identical setup as Section~4.1. We identify
index variables via AST analysis: loop counter variables bound in \texttt{for}
headers and dictionary key variables used in subscript expressions.

The baseline probe achieves F1~$= 0.959$ at L8 for Qwen2.5-1.5B and
F1~$= 0.945$ at L11 for DeepSeek-Coder-1.3B
(Figure~\ref{fig:appendix-sree-index-probe}), both higher than the accumulator
role ($0.893$) due to the strong lexical regularity of index names (\texttt{i},
\texttt{j}, \texttt{idx}). To isolate this lexical dependence, we apply
misleading renaming---swapping accumulator-style names onto index variables and
vice versa. As shown in Table~\ref{tab:index-renaming} and
Figure~\ref{fig:appendix-sree-index-renamed}, both models suffer the largest F1
drop of any role we study: $\Delta\text{F1} = {-}0.272$ for Qwen2.5-1.5B and
$\Delta\text{F1} = {-}0.285$ for DeepSeek-Coder-1.3B. This contrasts sharply
with the accumulator role, where the same manipulation yields a \emph{positive}
$\Delta$ F1. Cross-language transfer is strong across both models and exceeds
the accumulator range (Figure~\ref{fig:appendix-sree-index-crosslang}),
consistent with the syntactic regularity of loop-index patterns across languages.

\begin{table}[t]
\centering
\caption{Index/key probe F1 under key naming conditions across models, and
comparison with the accumulator role (Qwen2.5-1.5B). Both models show the
largest drop under misleading renaming, opposite to the accumulator pattern.}
\label{tab:index-renaming}
\begin{tabular}{llccc}
\toprule
Condition & Model & Best Layer & Test F1 & $\Delta$ F1 \\
\midrule
\multicolumn{5}{l}{\textit{Index role}} \\
Baseline          & Qwen2.5-1.5B        & L8  & 0.959 & ---      \\
Baseline          & DeepSeek-Coder-1.3B & L11 & 0.945 & ---      \\
All same (\texttt{x}) & Qwen2.5-1.5B   & L14 & 0.742 & $-$0.217 \\
Misleading        & Qwen2.5-1.5B        & L15 & 0.687 & $-$0.272 \\
Misleading        & DeepSeek-Coder-1.3B & L16 & 0.660 & $-$0.285 \\
\midrule
\multicolumn{5}{l}{\textit{Accumulator role (Qwen2.5-1.5B), for contrast}} \\
Baseline          & Qwen2.5-1.5B        & L9  & 0.893 & ---      \\
Misleading        & Qwen2.5-1.5B        & L9  & 0.972 & $+$0.079 \\
\bottomrule
\end{tabular}
\end{table}

Together, the accumulator and index results reveal a fundamental dissociation:
the model tracks index roles largely through lexical identity, while it tracks
accumulator roles through structural and positional context. This pattern holds
across both Qwen2.5-1.5B and DeepSeek-Coder-1.3B.
